% paper/sections/12f_error_budget.tex
%
% §12.6  精度バジェット（Error Budget）
%   §12.1--§12.5 の物理的整合性検証を踏まえた精度設計の総括
%   物理ベンチマーク（静止液滴・毛細管波・気泡上昇・RT不安定）は
%   第~\ref{sec:validation}章（§12）に収録．

% ═══════════════════════════════════════════════════════════════════════════════
\subsection{精度バジェット}
\label{sec:error_budget}
% ═══════════════════════════════════════════════════════════════════════════════

本章の検証結果を統合し，各コンポーネントの実測精度を一覧化する．
物理的な多相流ベンチマーク（静止液滴・毛細管波・気泡上昇）は §\ref{sec:validation} で扱う．

\begin{table}[htbp]
\centering
\caption{コンポーネント精度一覧（実測値）}
\label{tab:error_budget}
\renewcommand{\arraystretch}{1.35}
\begin{tabular}{>{\raggedright\arraybackslash}p{32mm}>{\raggedright\arraybackslash}p{22mm}>{\raggedright\arraybackslash}p{22mm}>{\raggedright\arraybackslash}p{16mm}>{\raggedright\arraybackslash}p{38mm}}
\toprule
コンポーネント & 設計精度 & 実測精度 & 検証箇所 & 誤差が影響する物理量 \\
\midrule
CCD 微分（内部） & $\Ord{h^6}$ & $\Ord{h^{5.97}}$ & §~\ref{sec:verify_ccd_convergence} & 速度微分，PPE 右辺 \\
CCD-PPE（変密度） & $\Ord{h^6}$ & $\Ord{h^{5.8}}$ & §~\ref{sec:verify_ppe_varrho} & 圧力場，速度補正 \\
CCD-PPE（$\varrho{=}2$） & --- & $\Ord{h^{1\text{--}2}}$ & §~\ref{sec:varrho_grid_convergence} & 変密度圧力場 \\
HFE 2D 場延長 & $\Ord{h^6}$ & $\Ord{h^{3.0}}$ & §~\ref{sec:verify_field_extension} & IPC $\bnabla p^n$，速度補正 \\
DCCD CLS 移流 & $\Ord{\varepsilon_d h^2}$ & $\Ord{h^{2}}$\textsuperscript{$\dagger$} & §~\ref{sec:verify_cls_advection} & 界面位置，体積保存 \\
CCD 曲率 $\kappa$ & $\Ord{h^6}$ & $\Ord{h^{5.0}}$ & §~\ref{sec:verify_curvature} & 表面張力体積力 \\
AB2+IPC 時間（単相） & $\Ord{\Delta t^2}$ & $\Ord{\Delta t^{2.00}}$ & §~\ref{sec:tgv_temporal} & 全体時間精度 \\
NS 空間（Kovasznay） & $\Ord{h^4}$--$\Ord{h^6}$ & $\Ord{h^{3.96}}$ & §~\ref{sec:kovasznay} & NS 運動量残差 \\
CSF 表面張力 & $\Ord{\varepsilon^2}$ & $\Ord{h^{1.5\text{--}2}}$ & §~\ref{sec:bf_static_droplet} & Laplace 圧精度 \\
\midrule
\multicolumn{5}{l}{\textit{追加検証}} \\
\midrule
混合偏微分 $\partial_{xy}$ & $\Ord{h^6}$ & $\Ord{h^{6.0}}$（周期） & §~\ref{sec:verify_mixed_partial} & 交差粘性項 \\
CN 粘性項 & $\Ord{\Delta t^2}$ & $\Ord{\Delta t^{2.0}}$ & §~\ref{sec:verify_cn_viscous} & 粘性時間積分 \\
DGR 厚み制御 & $\varepsilon_{\mathrm{eff}}/\varepsilon \approx 1$ & 1.03 & §~\ref{sec:verify_dgr} & 界面厚維持 \\
CCD vs FD 寄生流れ & CCD $\ll$ FD & FD/CCD $= 11\times$ & §~\ref{sec:force_balance} & Balanced-Force 効果 \\
交差粘性 CFL & $\Delta t \propto h^2/\Delta\mu$ & $C_{\mathrm{cross}} \approx 0.23$ & §~\ref{sec:cross_viscous_cfl} & 安定性制約 \\
毛細管 CFL & $\Delta t_\sigma \propto \Delta x^{3/2}$ & 指数 $= 1.505$ & §~\ref{sec:advection_pressure_coupling} & 時間ステップ制限 \\
\bottomrule
\end{tabular}
\end{table}

% ─────────────────────────────────────────────────────────────────────────────
\paragraph{律速段階の特定}
\label{sec:accuracy_summary}

\noindent\textbf{単相 NS ソルバ：}
空間精度は CCD 境界スキーム（$\Ord{h^4}$--$\Ord{h^5}$）が
非線形対流項で増幅され $\Ord{h^{3.96}}$（§~\ref{sec:kovasznay}）．
時間精度は AB2+IPC の $\Ord{\Delta t^2}$ が律速
（van Kan 1986）．
各コンポーネントの設計精度は NS パイプライン全体でほぼ完全に発現しており，
単相ソルバとしての精度は設計通りである．

\noindent\textbf{二相流（$\rho_l/\rho_g \leq 5$，smoothed Heaviside）：}
空間精度は CSF モデル誤差 $\Ord{\varepsilon^2} \approx \Ord{h^2}$ が律速する
（表~\ref{tab:bf_laplace}）．
CCD-PPE（$\Ord{h^6}$）は CSF より高次であり，
変密度一括 PPE は界面近傍で $\Ord{h^{1\text{--}2}}$ に劣化するが
CSF と同次である（§~\ref{sec:interface_crossing}）．
静止界面は 200 ステップ以上安定に動作する
（表~\ref{tab:static_longterm}）が，
移動界面＋表面張力では IPC の $\bnabla p^n$ 評価が界面ジャンプで
破綻する（表~\ref{tab:galilean}）．

\noindent\textbf{PPE 離散化：}
CCD PPE の欠陥補正は $L_L^{-1} L_H$ の不定性により
定常反復法では収束不能であり（§~\ref{sec:ppe_discretization_choice}），
FD 直接法 + CCD 勾配を採用する．
投影不整合 $\Ord{h^2}$ は有界であり（§~\ref{sec:interface_crossing}），
CSF モデル誤差 $\Ord{h^2}$ と同次のため全体精度を律速しない．

\noindent\textbf{二相流（$\rho_l/\rho_g \geq 10$）：}
変密度一括 PPE の欠陥補正は発散する（表~\ref{tab:varrho_breakdown}）が，
分相 PPE + DC $k{=}3$ により $\Ord{h^{7.0}}$ で全密度比をカバーする
（§~\ref{sec:verify_split_ppe_dc}）．

% ─────────────────────────────────────────────────────────────────────────────
\paragraph{CCD 高次精度の意義}

CCD の $\Ord{h^6}$ 精度は現状の CSF 近似（$\Ord{h^2}$）のもとでは
\textbf{直接的な全体精度向上には寄与しない}．
しかし以下の 2 つの経路で本質的な役割を果たす：
\begin{enumerate}[noitemsep]
  \item \textbf{寄生流れの係数低減：}
        Balanced-Force 条件下では収束\textbf{次数}は CSF 律速で $\Ord{h^2}$ だが，
        収束\textbf{係数}は曲率精度に比例する．
        CCD（$p_\kappa = 6$）は 2 次法と比べ $h = 1/64$ で
        $(1/64)^4 \approx 10^{-7}$ 倍の係数低減効果を持つ
        （§~\ref{sec:balanced_force}）．
  \item \textbf{分相 PPE 解法による精度回復：}
        分相 PPE 解法は CSF の $\Ord{h^2}$ モデル誤差を原理的に除去し，
        各相の PPE を独立に解くことで
        CCD の $\Ord{h^6}$ 設計精度を回復する
        （§~\ref{sec:interface_crossing}）．
\end{enumerate}

% ─────────────────────────────────────────────────────────────────────────────
\paragraph{ソルバの適用範囲}

表~\ref{tab:usage_constraints} に，本章の検証結果から導かれるソルバの適用制約を集約する．

\begin{table}[htbp]
\centering
\caption{現行ソルバの適用制約（本章の検証に基づく）}
\label{tab:usage_constraints}
\renewcommand{\arraystretch}{1.3}
\small
\begin{tabular}{p{42mm}ccp{38mm}}
\toprule
条件 & 対応 & 参照 & 備考 \\
\midrule
単相 NS（任意 Re） & \checkmark & §~\ref{sec:ns_time_accuracy} & $\Ord{\Delta t^2}$，$\Ord{h^4}$ 実効精度 \\
静止界面＋$\sigma > 0$，$\rho_l/\rho_g \leq 5$ & \checkmark & §~\ref{sec:twophase_static_longterm} & 200 ステップ安定 \\
移動界面，$\sigma = 0$，$\rho_l/\rho_g \leq 5$ & \checkmark & §~\ref{sec:galilean_invariance} & Galilean 不変，RT 誤差 2.8\% \\
移動界面＋$\sigma > 0$（任意 $\rho$） & $\times$ & §~\ref{sec:coupling_limitations} & IPC $\bnabla p^n$ が界面跳びで発散 \\
$\rho_l/\rho_g \geq 10$（smoothed H） & $\times$ & §~\ref{sec:interface_crossing} & DC 発散；分相 PPE で解決 \\
分相 PPE + DC $k{=}3$ & \checkmark & §~\ref{sec:verify_split_ppe_dc} & 全密度比で $\Ord{h^{7.0}}$ \\
\bottomrule
\end{tabular}
\end{table}

% ─────────────────────────────────────────────────────────────────────────────
\paragraph{§~\ref{sec:validation}章への接続}

§~\ref{sec:twophase_static_longterm} で低密度比の静止界面二相流が安定動作することを確認し，
§~\ref{sec:galilean_invariance} で移動界面＋表面張力には圧力場の分相延長が必要であることを
特定した．
§~\ref{sec:interface_crossing} で smoothed Heaviside の破綻限界を定量化し，
§~\ref{sec:verify_split_ppe_dc} では分相 PPE + DC $k{=}3$ により
$\rho_l/\rho_g = 1000$ でも $\Ord{h^{7.0}}$ 格子収束が達成されることを実証した
（表~\ref{tab:split_ppe_dc_convergence}）．
第~\ref{sec:validation}章では現行ソルバの動作範囲内（$\rho_l/\rho_g \leq 5$，
静止界面＋CSF）で RT 不安定と静止液滴の物理ベンチマークを実施し，
HFE（§~\ref{sec:field_extension}）の除去実験と
密度比スイープにより本手法の能力と限界を定量化する．

\noindent\textbf{検証対象外の保存量：}
角運動量はコロケート格子上の離散 NS ソルバでは一般に厳密保存されない%
\footnote{スタガード格子は運動量の局所保存を格子構造により保証するが，
コロケート格子では圧力--速度の補間構造上，
角運動量フラックスの厳密な保存は保証されない．
本稿では DCCD フィルタ（$\varepsilon_d = 1/4$）による欠陥補正で
チェッカーボード不安定性を抑制しており（§\ref{sec:dccd_decoupling}），
面流速補間に起因する角運動量誤差は副次的である．
角運動量保存の厳密な検証にはスタガード格子への拡張が必要であり，
本稿の範囲外とする．}．
したがって本章では角運動量保存の独立検証は実施しない．
